\documentclass[a4paper, serif, 10pt]{moderncv}

%% moderncv
% style and color
\moderncvstyle{banking}
\moderncvcolor{black}

% renew moderncv command to adapt for CJK colon
\renewcommand*{\cvitem}[3][.25em]{%
  \ifstrempty{#2}{}{\hintstyle{#2}：}{#3}%
  \par\addvspace{#1}}

\renewcommand*{\cvitemwithcomment}[4][.25em]{%
  \savebox{\cvitemwithcommentmainbox}{\ifstrempty{#2}{}{\hintstyle{#2}：}#3}%
  \setlength{\cvitemwithcommentmainlength}{\widthof{\usebox{\cvitemwithcommentmainbox}}}%
  \setlength{\cvitemwithcommentcommentlength}{\maincolumnwidth-\separatorcolumnwidth-\cvitemwithcommentmainlength}%
  \begin{minipage}[t]{\cvitemwithcommentmainlength}\usebox{\cvitemwithcommentmainbox}\end{minipage}%
  \hfill% fill of \separatorcolumnwidth
  \begin{minipage}[t]{\cvitemwithcommentcommentlength}\raggedleft\small\itshape#4\end{minipage}%
  \par\addvspace{#1}}

%% page layout/margins
\usepackage[top=2.5cm, bottom=2.5cm, left=1.5cm, right=1.5cm]{geometry}

\nopagenumbers{}

%% Babel config for English language
\usepackage[english]{babel}

%% fontspec
\usepackage{fontspec}

\IfFontExistsTF{Linux Libertine}{
  \setmainfont[Ligatures={TeX, Common}, Numbers=Lining, ItalicFont=Linux Libertine]{Linux Libertine}
}{}
\IfFontExistsTF{Linux Libertine O}{
  \setmainfont[Ligatures={TeX, Common}, Numbers=Lining, ItalicFont=Linux Libertine O]{Linux Libertine O}
}{}

%% CTeX
% CJK support, used to show CJK characters in the resume
%
% - fontset=none: disable builtin fontset but instead set the CJK font manually
% - heading=false: disable ctex heading
% - punct=kaiming: use kaiming punctuations styles for CJK
% - scheme=plain: use plain scheme, do not override \normalsize font size
% - space=auto: space settings for CJK characters
%
% ref:
% - http://ctan.mirrorcatalogs.com/language/chinese/ctex/ctex.pdf
\usepackage[UTF8, heading=false, punct=kaiming, scheme=plain, space=auto]{ctex}

\IfFontExistsTF{Noto Serif CJK SC}{
  \setCJKmainfont{Noto Serif CJK SC}
}{}
\IfFontExistsTF{Noto Sans CJK SC}{
  \setCJKsansfont{Noto Sans CJK SC}
}{}

%% line spacing
\usepackage{setspace}
\setstretch{1.125}

%% URL styling - use normal text instead of monospace
\urlstyle{same}

% Auto-underline all links
% Defer until \begin{document} so hyperref is already loaded
\AtBeginDocument{%
  \let\oldhref\href
  \renewcommand{\href}[2]{\oldhref{#1}{\underline{#2}}}%
}

%% Patch \httplink and \httpslink to handle full URLs
% The original moderncv macros always prepend http:// or https:// to URLs.
% This causes issues when the URL already has a protocol (e.g., https://example.com)
% resulting in malformed URLs like https://https://example.com.
% This patch checks if the URL already contains "://" and uses it directly if so.
% Note: We use \str_set:Nx (x-expansion) to expand macros like \@homepage first.
\ExplSyntaxOn
\str_new:N \g_yamlresume_url_str

\RenewDocumentCommand{\httpslink}{O{}m}{
  \str_set:Nx \g_yamlresume_url_str {#2}
  \str_if_in:NnTF \g_yamlresume_url_str {://}
    { \url{#2} }
    { \url{https://#2} }
}

\RenewDocumentCommand{\httplink}{O{}m}{
  \str_set:Nx \g_yamlresume_url_str {#2}
  \str_if_in:NnTF \g_yamlresume_url_str {://}
    { \url{#2} }
    { \url{http://#2} }
}
\ExplSyntaxOff

\name{林思远}{}
\title{客户成功经理 | SaaS 客户增长与续约专家}
\phone[mobile]{+86 138 0013 8000}
\email{siyuan.lin@example.com}
\homepage{https://siyuanlin.example.com}

\address{中国，北京市，北京，北京市朝阳区建国路88号，100022}{}{}

\extrainfo{{\small \faLinkedin}\ \href{https://www.linkedin.com/in/siyuanlin-csm}{@siyuanlin-csm} {} {} {} • {} {} {} 
{\small \faWeixin}\ \href{https://wechat.com/siyuanlin\_csm}{@siyuanlin\_csm} {} {} {} • {} {} {} 
{\small \faZhihu}\ \href{https://www.zhihu.com/people/siyuanlin-csm}{@siyuanlin-csm}}

\begin{document}

\maketitle

\section{简介}

\cvline{}{深耕 SaaS 与科技行业 7 年，专注客户成功、续约增购与客户体验管理。擅长通过数据驱动的客户健康度模型识别风险与机会，推动跨部门协作，提升客户生命周期价值。
%
\begin{itemize}
\item 负责 30+ 家中大型企业客户，管理年合同金额超 3000 万元，续约率连续两年保持在 95\% 以上
\item 主导客户 onboarding、QBR 与成功计划，客户 NPS 从 32 提升至 68
\item 熟练使用 Salesforce、Gainsight、Jira 等工具，具备优秀的中英文沟通与谈判能力
\item 以客户价值为导向，善于将客户反馈转化为产品改进与增购机会
\end{itemize}}
\section{教育背景}

\cventry{2012年9月–2016年6月}
        {学士，工商管理，成绩：3.7/4.0}
        {上海交通大学}
        {\url{https://www.sjtu.edu.cn}}
        {}
        {\begin{itemize}
\item 系统学习管理学、市场营销与数据分析，为客户成功与商业洞察打下坚实基础
\item 担任商学院职业发展协会外联负责人，组织 10+ 场企业参访与求职分享活动
\item 获校级优秀学生奖学金两次，GPA 3.7/4.0
\end{itemize}
\textbf{课程}：管理学原理、市场营销、组织行为学、消费者行为、商务统计分析、运营管理、企业战略管理、商务沟通}
\section{工作经历}

\cventry{2021年7月至今}
        {高级客户成功经理}
        {云启科技}
        {\url{https://www.yunqi-tech.com}}
        {}
        {\begin{itemize}
\item 负责 30+ 家中大型企业客户，涵盖制造、零售与互联网行业，管理年合同金额超 3000 万元
\item 建立客户健康度评分体系，结合产品使用率、工单趋势与关键联系人活跃度，提前识别续约风险
\item 主导季度业务回顾（QBR）与成功计划，推动续约率连续两年保持在 95\% 以上
\item 挖掘客户增购与交叉销售机会，年度增购收入贡献超过 800 万元
\item 与产品、研发、销售团队紧密协作，将高频客户反馈转化为产品需求，推动 5 项核心功能上线
\item 带领 3 名客户成功专员，制定 onboarding 标准流程与客户培训体系
\end{itemize}
\textbf{关键字}：SaaS客户成功、续约管理、客户增购、NPS提升、跨部门协作}

\cventry{2018年8月–2021年6月}
        {客户成功专员}
        {智联数据}
        {\url{https://www.zhilian-data.com}}
        {}
        {\begin{itemize}
\item 负责 60+ 家中小型企业客户的日常运营与关系维护，客户满意度保持在 90\% 以上
\item 搭建客户培训与知识库体系，制作 20+ 份产品操作指南与最佳实践文档
\item 通过定期回访与使用数据分析，成功挽回 15 家高风险续约客户
\item 协同销售团队完成新客户交接，平均缩短 onboarding 周期 30\%
\end{itemize}
\textbf{关键字}：客户关系维护、客户培训、风险预警、数据分析}

\cventry{2016年7月–2018年7月}
        {客户支持顾问}
        {恒信咨询}
        {\url{https://www.hengxin-consulting.com}}
        {}
        {\begin{itemize}
\item 为金融与零售客户提供系统上线支持与业务咨询服务，累计服务 100+ 家企业用户
\item 处理客户疑难问题与投诉，协调内部资源，平均问题解决时长缩短至 4 小时以内
\item 参与客户需求调研与方案撰写，协助销售团队完成 20+ 个项目的售前支持
\end{itemize}
\textbf{关键字}：客户支持、售前咨询、问题解决、CRM}
\section{语言}

\cvline{中文}{母语或双语水平 \hfill \textbf{关键字}：普通话、商务写作}
\cvline{英语}{完全专业水平 \hfill \textbf{关键字}：CET-6、商务谈判、英文演示}
\section{专业技能}

\cvline{客户成功管理}{专家 \hfill \textbf{关键字}：客户生命周期管理、续约与增购、QBR与成功计划、客户健康度模型、onboarding流程}
\cvline{数据分析与洞察}{高级 \hfill \textbf{关键字}：SQL、Excel、Tableau、客户分群、流失预警}
\cvline{CRM与协作工具}{高级 \hfill \textbf{关键字}：Salesforce、Gainsight、Jira、Confluence、企业微信}
\cvline{沟通与谈判}{专家 \hfill \textbf{关键字}：商务谈判、高层沟通、冲突处理、中英文演示}
\cvline{项目管理}{中级 \hfill \textbf{关键字}：跨部门协作、敏捷方法、需求管理、项目排期}
\section{荣誉}

\cventry{2023年12月}
        {云启科技}
        {年度客户成功之星}
        {}
        {}
        {因主导客户健康度体系与续约专项，年度续约率与增购收入均超额完成目标，获公司年度客户成功团队最高个人荣誉。}

\cventry{2020年1月}
        {智联数据}
        {优秀员工}
        {}
        {}
        {连续两个季度客户满意度排名第一，并成功挽回多家高价值客户续约。}
\section{证书}

\cventry{2020年5月}
        {客户成功协会}
        {客户成功经理认证（CCSM）}
        {\url{https://www.customersuccessassociation.com/certification}}
        {}
        {}

\cventry{2019年9月}
        {项目管理协会（PMI）}
        {项目管理专业人士（PMP）}
        {\url{https://www.pmi.org/certifications/project-management-pmp}}
        {}
        {}
\section{出版刊物}

\cventry{2022年4月}
        {SaaS 客户健康度模型实践指南}
        {人人都是产品经理}
        {\url{https://www.woshipm.com/practice/example.html}}
        {}
        {系统介绍如何从产品使用、服务交互、关键联系人活跃度与商业价值四个维度构建客户健康度评分，并结合实际案例说明如何驱动续约与增购。}
\section{项目}

\cventry{2022年3月–2022年10月}
        {面向 SaaS 客户成功团队的客户风险识别与增购机会挖掘系统。}
        {客户健康度评分体系}
        {\url{https://www.yunqi-tech.com/customer-health}}
        {}
        {\begin{itemize}
\item 设计 4 大维度、18 项指标的客户健康度模型，覆盖产品使用、服务体验、关系强度与商业价值
\item 联合数据团队搭建自动化看板，每日更新客户评分与预警名单
\item 上线后帮助团队提前 60 天识别 80\% 的续约风险客户，续约率提升 8\%
\end{itemize}
\textbf{关键字}：客户健康度、数据看板、续约预警、SaaS}

\cventry{2019年5月–2020年12月}
        {为中小企业客户设计的线上培训与知识库赋能项目。}
        {客户培训与赋能计划}
        {\url{https://www.zhilian-data.com/customer-training}}
        {}
        {\begin{itemize}
\item 梳理客户高频问题与使用场景，制作 20+ 份操作指南与视频课程
\item 搭建自助知识库，客户自助解决率从 35\% 提升至 65\%
\item 组织月度线上培训，累计参与客户超过 500 人次
\end{itemize}
\textbf{关键字}：客户培训、知识库、自助服务}
\section{兴趣爱好}

\cvline{阅读}{商业管理、心理学、科技趋势}
\cvline{运动}{徒步、羽毛球、游泳}
\section{志愿服务}

\cventry{2019年3月–2021年12月}
        {职业发展导师}
        {北京市朝阳区社区服务中心}
        {\url{https://www.example-community.org}}
        {}
        {\begin{itemize}
\item 为社区青年提供职业规划与求职辅导，累计服务 200+ 人次
\item 组织 SaaS 行业分享会，帮助求职者了解客户成功岗位技能要求
\item 协助社区开展公益培训活动，获评年度优秀志愿者
\end{itemize}}
    
\end{document}